% packages needed in all documents

\usepackage[T1]{fontenc}
\usepackage[utf8]{inputenc}
\usepackage{ngerman}
\usepackage{verbatim}
\usepackage[hyphens,spaces]{url}
\usepackage{listings}
\usepackage{lastpage}
\usepackage{graphicx}
\usepackage{ifthen}
\usepackage{index}
\usepackage{color}
\usepackage{placeins}
\usepackage{float}
\usepackage{longtable}
\usepackage{pifont}
\usepackage{version}
\usepackage{multirow}
\usepackage{eso-pic}
\usepackage{wallpaper}
\usepackage[top=4.5cm, headsep=1.5cm, bottom=3cm, footskip=2cm]{geometry}
%\usepackage{showkeys}  % ENABLE SHOWKEYS TO DISPLAY REFERENCE_TAGS IN DOCUMENT
\usepackage{ucsframed}
\usepackage{marginnote}
\usepackage{placeins}

\def\UrlBreaks{\do\/\do\.\do\-}

\iftrue
% http://www.ctan.org/tex-archive/macros/latex/required/psnfss/psnfss2e.pdf
\usepackage{mathptmx} % roman, formulas
\usepackage{helvet} % sans-serif
\usepackage{courier} % typewriter
\renewcommand{\familydefault}{\sfdefault}
\else
% use font univers as sf-font
\fontfamily{pun}
\selectfont
\sffamily
\renewcommand*{\rmdefault}{pun} % pageheader
\renewcommand*{\sfdefault}{pun} % global font
%\renewcommand*{\ttdefault}{pun} % fixed font, should not be univers
\fi

% Tabbox, can be used to create s.th. like this
%
% ---------------------
% | Tab-Titel         |
% ------------------------------------
% | whatever you want                |
% | and some more of this            |
% ------------------------------------
%
% usage:
%
% \begin{ucsTabbox}{Tab-Titel}
% whatever you want
% and some more of this
% \end{ucsTabbox}

% Tabbox with Info Sign

\newenvironment{ucsInfobox}[1]
{\marginpar{\vspace{3ex}\centering\includegraphics{abbildungen/info.jpg}}\setlength{\FrameRule}{1pt}\setlength{\fboxsep}{2mm}\begin{ucsframed}{#1}\parskip\baselineskip}
{\end{ucsframed}}

\newenvironment{ucsTabbox}[1]
{\setlength{\FrameRule}{1pt}\setlength{\fboxsep}{2mm}\begin{ucsframed}{#1}\parskip\baselineskip}
{\end{ucsframed}}

\newcommand{\ucsTabentry}[2]{\textbf{#1}\\#2}

%% for the background-image on the title page
%% http://stackoverflow.com/questions/240097/how-to-create-a-background-image-on-titlepage-with-latex
\newcommand\BackgroundPicEven{
\put(0,0){
\parbox[b][\paperheight]{\paperwidth}{%
\vfill
\centering
\includegraphics[width=\paperwidth,height=\paperheight]{illustrations/page-background-even.jpg}%
\vfill
}}}

\newcommand\BackgroundPicOdd{
\put(0,0){
\parbox[b][\paperheight]{\paperwidth}{%
\vfill
\centering
\includegraphics[width=\paperwidth,height=\paperheight]{illustrations/page-background-odd.jpg}%
\vfill
}}}

% Colors:
% red as used in univention corporate identiy
% gray for listing background

\definecolor{ucsRed}{RGB}{0.51372549, 0.141176471, 0.235294118}
\definecolor{ucsCodeGray}{gray}{0.7}
\definecolor{ucsHyperrefLinkColor}{rgb}{0, 0, 0.8}

% tabular environment using sans serif

\newenvironment{tableSansSerif}%
{\begin{table}[H]\footnotesize\sffamily}%
{\end{table}}

% environment for console sessions
\lstnewenvironment{ucsConsoleInput}
 {\lstset{basicstyle=\ttfamily}\small}
 {\normalsize}

% environment for console sessions, boxed
%\lstnewenvironment{ucsConsoleInputBoxed}
% {\lstset{frame=trlb,basicstyle=\ttfamily}\small}
% {\normalsize}

% environment for config files
\lstnewenvironment{ucsConfigFile}
 {\lstset{basicstyle=\ttfamily,backgroundcolor=\color{ucsCodeGray}}\small}
 {\normalsize}

% exclamation mark on border
\newcommand{\ucsExclamation}{\mbox{}\marginnote{\vspace{0pt}\centering\includegraphics{abbildungen/achtung.jpg}}}

% URLs
\newcommand{\ucsURL}[1]{\textsf{\mbox{\url{#1}}}}
%\newcommand{\ucsURL}[1]{\textsf{\url{#1}}}

\newcommand{\ucsFile}[1]{\sloppy{\texttt{#1}}}

% Bugs
\newcommand{\ucsBug}[2][Bug \#]{\href{https://forge.univention.org/bugzilla/show_bug.cgi?id=#2}{#1#2}}

% emphasize text
\renewcommand{\emph}[1]{\textbf{\textsl{#1}}}

% various names
\newcommand{\ucsName}[1]{\sloppy{\textbf{\textsl{#1}}}}

% single command, inline
\newcommand{\ucsCommand}[1]{\sloppy{\texttt{#1}}}

% easy backslash
\newcommand{\BS}{\ensuremath{\backslash}}

% menu entries
\newcommand{\ucsMenuEntry}[1]{\textbf{#1}}

% click buttons
\newcommand{\ucsButton}[1]{\sloppy{[\textbf{#1}]}}

% special notes
\newcommand{\ucsNote}[2][Hinweis:]{\textsl{\textbf{#1}\ifthenelse{\equal{#1}{}}{}{\\}#2}}

% extra note
\newcommand{\ucsAttention}[1]{\textbf{Achtung:}\\#1}

% example
\newcommand{\ucsExample}[2][Beispiel:]{\textbf{#1}\\#2}

% nice right arrow
\newcommand{\ucsRightArrow}{\ding{222}}

% no new paragraph indention
\parindent0cm

% easy integration of bitmaps
\newcommand{\ucsGraphicsTabbox}[3][1]
 {
  \begin{center}
   \includegraphics[width=#1\textwidth]{#3}
   {\refstepcounter{figure}Abbildung \thefigure: \begin{minipage}[t]{0.7\textwidth}{#2}\end{minipage}}
  \end{center}
 }

% easy integration of bitmaps
\newcommand{\ucsGraphics}[3][1]
 {
  \begin{figure}[htb]
  \begin{center}
   \includegraphics[width=#1\textwidth]{#3}
   \ifthenelse{\equal{#2}{}}{}{\caption{#2}}
  \end{center}
  \end{figure}
 }

% easy integration of bitmaps (with reference)
\newcommand{\ucsGraphicsRef}[4][1]
 {
  \begin{figure}[htb]
  \begin{center}
   \includegraphics[width=#1\textwidth]{#3}
   \ifthenelse{\equal{#2}{}}{}{\caption{#2}}
   \label{pic:#4}
  \end{center}
  \end{figure}
 }

% simple index command for univention baseconfig variables
%
% PLEASE do not insert additional spaces or newlines for readability
% they will occur as additional and annoying spaces within final document!
\newcommand{\ucrpath}[1]{% <http://tex.stackexchange.com/questions/50777/>
 \begingroup
 \ttfamily
 \begingroup\lccode`~=`/\lowercase{\endgroup\def~}{/\discretionary{}{}{}}%
 \catcode`/=\active
 \scantokens{#1\noexpand}%
 \endgroup
}
\newcommand{\ucsBCindex}[2][]{\ifthenelse{\equal{#1}{}}{#2}{#1}\index[ucsConfigRegistry]{#2@\ucrpath{#2}}}

\tracingpages=1

% add section number up to subparagraph
\setcounter{secnumdepth}{5}
% add section into index up to subsection
\setcounter{tocdepth}{5}

%wording
\newcommand{\ucsUDM}[1]{Univention Directory Manager#1}
\newcommand{\ucsUDL}[1]{Univention Directory Listener#1}
\newcommand{\ucsUCD}[1]{Univention Corporate Desktop#1}
\newcommand{\ucsUDN}[1]{Univention Directory Notifier#1}
\newcommand{\ucsUMS}[1]{UCS-Managementsystem#1}
\newcommand{\ucsTCS}[1]{Univention Thin Client Services#1}
\newcommand{\ucsUMC}[1]{Univention Management Console#1}
\newcommand{\ucsUCR}[1]{Univention Configuration Registry#1}
\newcommand{\ucsUSS}[1]{Univention System Setup#1}
\newcommand{\ucsUCRV}[1]{\ucsUCR variable \ucsCommand{\ucsBCindex{#1}}}
\newcommand{\ucsUCRVSA}[1]{\ucsCommand{\ucsBCindex{#1}}}
\newcommand{\ucsUAS}[1]{UCS@school#1}
\newcommand{\ucsUVMM}[1]{UCS Virtual Machine Manager#1}

\newcommand{\ucsDVS}[1]{UCS Desktop Virtualization Services#1}
\newcommand{\ucsMaster}[1]{Domänencontroller Master#1}
\newcommand{\ucsBackup}[1]{Domänencontroller Backup#1}
\newcommand{\ucsSlave}[1]{Domänencontroller Slave#1}
\newcommand{\ucsMember}[1]{Memberserver#1}

\newcommand{\ucsUCS}[1]{Univention Corporate Server#1}

\newcommand{\ucsUGS}[1]{Univention Groupware Server#1}
\newcommand{\ucsKolab}[1]{Kolab2 für UCS#1}

%some \ucsUDMb,\ucsUDLCb and \ucsUCRb entrys have to be changed by hand (they do not work in listing environments)
%like  "\begin{ucsConsoleInput}" or something like that
\newcommand{\ucsUCRb}[1]{univention-config-registry#1}
\newcommand{\ucsUDMb}[1]{univention-directory-manager#1}
\newcommand{\ucsUDLCb}[1]{univention-directory-listener-ctrl#1}
\newcommand{\ucsUDNb}[1]{univention-directory-notifier#1}

%not bash but bash style, because there is no command "univention-management-console"
\newcommand{\ucsUMCnb}[1]{univention-management-console#1}

